\section{The Twin, Re-Scored, and Two Corrections Against Ourselves}
\label{sec:twin}

An earlier version of this work reported that the dynamic twin scores CSI 0.041 against Sentinel-1,
that static terrain indices score the same, and concluded that \emph{no} terrain-based router can
place ponding in this data regime. Two of those three statements survive re-measurement. The third
does not, and we correct it here, because the experiment that refutes it is our own.

\subsection{Protocol}
Sentinel-1 water masks for 18 monsoon storms across two areas, Patna with 7 storms from 2023 to
2025 and Mumbai-Northeast with 11 storms from 2025 to 2026, are scored on a common 60\,m grid with
permanent water masked \emph{on both sides}. Every method, whether physical, index-based or
learned, sees the identical grid, mask and target. Every method's detection threshold is chosen on
the \emph{other} storms of the fold and never on the scene being reported. One Patna date,
2024-07-07, has zero observed wet cells and therefore no ground truth, and it contributed a forced
0.0 to every method's mean in our earlier table. It is dropped. We verified the scoring path by
re-deriving the earlier table from it under the earlier protocol, which moved the twin from 0.0412
to 0.0409, TWI from 0.0422 to 0.0429, and HAND-lite from 0.0505 to 0.0504.

We keep this like-for-like table restricted to these two areas deliberately. They are the only
areas for which a twin baseline exists to compare against, and mixing in areas where only the
learned model can be run would not be a fair head-to-head. The larger 15-area evidence base is used
everywhere else in this paper.

\begin{table}[!t]
\centering
\caption{All methods on the identical 18 storms, with every threshold chosen off the reported
scene. The learned model is a 7.9\,M-parameter U-Net over 23 terrain channels with rainfall
removed, averaged over 3 seeds. Every terrain method is quoted at its most favourable threshold
rule. POD is probability of detection, the fraction of truly wet cells found.}
\label{tab:baselines}
\begin{tabular}{lccc}
\toprule
method & mean CSI & bias & POD \\
\midrule
random                            & 0.014 & 1.00 & n/a \\
all-wet                           & 0.029 & n/a  & 1.00 \\
static depression depth           & 0.032 & n/a  & 0.21 \\
TWI (topographic wetness)         & 0.039 & n/a  & 0.26 \\
HAND-lite (height above drainage) & 0.039 & n/a  & 0.90 \\
dynamic twin (antecedent rain)    & 0.054 & n/a  & 0.29 \\
\midrule
\textbf{learned (terrain only)}   & \textbf{0.288 $\pm$ 0.007} & 1.41 & 0.52 \\
climatology (no model, no rain)   & 0.288 & 1.29 & n/a \\
\bottomrule
\end{tabular}
\end{table}

\subsection{Correction one: the twin is better than we published}
On a fair test set the twin reaches 0.054, not the 0.041 we reported. Dropping the dead scene and
adding a second city both help it. It also clears all-wet at 0.029 and every static index, which
our Patna-only table did not show. We state this plainly because the correction runs against our
own earlier headline in our favour, and a correction that only ever runs one way is not a
correction.

\subsection{Correction two: co-location is learnable, so our terrain-fundamental claim was wrong}
The same elevation, the same land cover and the same soil rasters, read by a U-Net instead of routed
by a solver, reach $0.288 \pm 0.007$. That is 5.4 times the twin and 9.9 times all-wet on identical
scenes. What fails in this regime is \emph{routing} water over a noisy elevation model, not
\emph{predicting} where water goes.

The feature that makes the difference is available to both methods. It is elevation minus its own
blurred copy at 2, 8 and 32 cells, which is a measure of whether a spot is locally low. Elevation
error is spatially correlated, meaning neighbouring cells are wrong together and by similar
amounts, so this relative micro-relief survives a bias that destroys absolute elevation.

\subsection{But the learned model ties the climatology it was meant to beat}
The learned model scores 0.2884 and the climatology scores 0.2881. That is the same number to three
decimal places, a gap of 0.0003 against a seed standard deviation of 0.0066. On a city that has a
Sentinel-1 archive, a neural network over 23 terrain channels is worth nothing over the statement
``these cells are usually wet''.

This is the hinge of the paper. The learned model's contribution is not accuracy. It is
portability, which we isolate in Section~\ref{sec:transfer}.

\subsection{An independent check with a different instrument}
Radar tells us where water was. It does not tell us how deep. To test depth we used crowdsourced
reports from the public mumbaiflood.in service, run by IIT Bombay's climate studies group with the
municipal corporation. Each report gives a reporter's height and the fraction of themselves the
water reached, so depth is the product of the two.

Of 147 raw reports, 38 carry no depth value, 15 fall outside the region because the app is
reachable worldwide and the live feed contains points in Argentina and Australia, and 2 report
more than a metre of standing water after under 4\,mm of rain, which is not rainfall-driven
waterlogging whatever else it is. That leaves 94, of which 85 fall inside a built Mumbai tile and
83 are scored. Personal fields are read and discarded at parse time, and a test pins that they
never reach a file. These observations train nothing. They are held for scoring only.

The result is bad for the twin, and it is bad in an informative way. Mean absolute error is
0.207\,m and root mean square error is 0.310\,m, with a bias of $-0.162$\,m, meaning systematic
under-prediction. Against the best trivial baseline, which is simply predicting the observed mean
depth everywhere, the skill score is $-0.46$. Negative skill means the model is \emph{worse} than
predicting one constant number for the whole city. No tile is positive, and the range runs from
$-0.40$ on Mumbai-West to $-0.64$ on Mumbai-South. Correlation with observed depth is $+0.055$
overall and $-0.23$ on the genuinely wet sites.

\subsection{Two failure modes, separated}
This experiment splits the failure cleanly into a part that is fixable and a part that is not.

\textbf{Magnitude is a resolution artifact, and it is curable.} On wet sites the model
under-predicts by a stable factor of 4.45, identical at the 0.15\,m and 0.30\,m thresholds. That
ratio is not mysterious. A person reports the depth where they are standing, and the model reports
the average over a 3\,600\,m$^2$ cell. If the same water sits in a pond of area $A$ inside that
cell, then observed over predicted is roughly cell area over $A$. That implies a ponding footprint
of $3\,600 / 4.45 \approx 810$\,m$^2$, a square 28.5\,m across. That is a flooded street junction.
The model is not wrong about how much water there is. It is spreading it over sixty metres of
ground.

\textbf{Ranking is the real failure, and resolution will not fix it.} Correlation with observed
depth is zero to negative, and it stays that way across every rainfall window and every
neighbourhood size we tried. The model cannot say which street floods worse. This is the same
co-location failure the radar work found, now confirmed by an independent data source measuring an
independent quantity. Two methods, two datasets, one verdict.

We state the limits of this result rather than leaning on it. There are 83 points from one city in
one season, so the negative correlation on wet sites is suggestive and not established. People
report water where it is surprising or bad, which is plausibly where the model least expects it,
and that reporting bias alone could produce anti-correlation without the model being wrong in the
way it appears to be. A point measurement and a cell average are genuinely different quantities,
so the 4.45 factor is a consistency argument rather than proof. And a report made hours after the
peak is compared against a storm maximum the model is entitled to predict.
